\documentclass[UTF8,12pt,a4paper,fontset=fandol]{ctexart}

\usepackage[a4paper,top=25mm,bottom=25mm,left=27mm,right=27mm]{geometry}
\usepackage{amsmath,amssymb,bm}
\usepackage{enumitem}
\usepackage{fancyhdr}
\usepackage{microtype}

\setlength{\parindent}{2em}
\setlength{\parskip}{0.25em}
\setlist[enumerate]{leftmargin=2.2em,itemsep=0.25em,topsep=0.35em}
\allowdisplaybreaks[2]

\pagestyle{fancy}
\setlength{\headheight}{15pt}
\fancyhf{}
\fancyhead[C]{二维对流--扩散方程数值测试}
\fancyfoot[C]{\thepage}
\renewcommand{\headrulewidth}{0.4pt}
\renewcommand{\footrulewidth}{0pt}
\makeatletter
\let\ps@plain\ps@fancy
\makeatother

\ctexset{
  section = {
    format = \Large\bfseries,
    number = \arabic{section},
    aftername = \quad
  },
  subsection = {
    format = \large\bfseries,
    number = \thesection.\arabic{subsection},
    aftername = \quad
  },
  subsubsection = {
    format = \normalsize\bfseries,
    number =,
    aftername =
  }
}

\begin{document}

\begin{center}
  {\LARGE\bfseries 第三题：二维对流--扩散方程}
\end{center}
\vspace{0.8em}

\setcounter{section}{2}
\section{对流--扩散方程}

\subsection{控制方程}

考虑二维对流--扩散方程

\begin{equation}
\frac{\partial\phi}{\partial t}
+\nabla\cdot(\bm{u}\phi)
-\nabla\cdot\left(\mu\nabla\phi\right)=0,
\end{equation}

其中，$\phi=\phi(x,y,t)$ 为待求标量，$\bm{u}$ 为给定速度场，$\mu$ 为扩散系数。

请基于有限体积法写出该方程的半离散形式，并分别采用显式时间离散格式和显式空间离散格式构建数值求解器。

\subsection{数值算例}

\subsubsection{3.2.1\quad 正弦波对流--扩散问题}

采用二维正弦波算例检验数值通量的精度。计算区域为

\begin{equation}
\Omega=[0,1]^2,
\end{equation}

初始条件为

\begin{equation}
\phi(x,y,0)=\sin\bigl(2\pi(x+y)\bigr).
\end{equation}

在 $x$、$y$ 两个方向施加周期边界条件：

\begin{equation}
\phi(0,y,t)=\phi(1,y,t),
\qquad
\phi(x,0,t)=\phi(x,1,t).
\end{equation}

给定均匀速度场和扩散系数

\begin{equation}
\bm{u}=(a,a),
\qquad
a=1,
\qquad
\mu=1.
\end{equation}

该初边值问题的精确解为

\begin{equation}
\phi_{\mathrm{ex}}(x,y,t)
=\exp\left(-8\pi^2\mu t\right)
\sin\bigl(2\pi(x+y-2at)\bigr).
\end{equation}

分别采用三角形单元和四边形单元对计算区域进行划分，并逐渐加密网格。

请在

\begin{equation}
t=1
\end{equation}

时计算数值误差和收敛阶，以评价数值解的精度与网格收敛行为。

\subsubsection{3.2.2\quad 旋转尖峰对流--扩散问题}

考虑计算区域

\begin{equation}
\Omega=[-1,1]^2
\end{equation}

上的非定常对流--扩散问题。给定旋转速度场和扩散系数

\begin{equation}
\bm{u}=(-y,x),
\qquad
\mu=\varepsilon=10^{-3}.
\end{equation}

定义随时间移动并扩散的尖峰精确解

\begin{equation}
\phi_{\mathrm{ex}}(x,y,t)
=\frac{1}{4\pi\varepsilon t}
\exp\left(-\frac{r^2(x,y,t)}{4\varepsilon t}\right),
\end{equation}

其中

\begin{equation}
r^2(x,y,t)
=\bigl(x-\widehat{x}(t)\bigr)^2
+\bigl(y-\widehat{y}(t)\bigr)^2,
\end{equation}

尖峰中心的运动轨迹为

\begin{equation}
\widehat{x}(t)=x_0\cos t-y_0\sin t,
\qquad
\widehat{y}(t)=-x_0\sin t+y_0\cos t,
\end{equation}

并取

\begin{equation}
x_0=0,
\qquad
y_0=0.5.
\end{equation}

从

\begin{equation}
t_0=\frac{\pi}{2}
\end{equation}

开始数值计算，以精确解在该时刻的取值作为初始条件：

\begin{equation}
\phi(x,y,t_0)=\phi_{\mathrm{ex}}(x,y,t_0),
\qquad (x,y)\in\Omega.
\end{equation}

在计算区域边界上施加由精确解给出的 Dirichlet 边界条件：

\begin{equation}
\phi(x,y,t)=\phi_{\mathrm{ex}}(x,y,t),
\qquad (x,y)\in\partial\Omega,
\quad t\geq t_0.
\end{equation}

速度场的角速度为 $1$，完成一周旋转所需的时间为 $2\pi$，因此终止时刻为

\begin{equation}
t_1=t_0+2\pi=\frac{5\pi}{2}.
\end{equation}

随着时间推进，尖峰将因扩散作用逐渐展宽。请绘制一周旋转后，即 $t=5\pi/2$ 时的标量场等值线图，并通过数值误差和网格收敛行为评价数值解的精度。

\end{document}
